\domainevide{Mathlib contient les sites et les faisceaux, mais pas les morphismes de topos, leurs points ni les topos classifiants, que ces cotes emploient.}{23, 37, 90, 112, 116, 156-6, 161-3, 161-4}
\noindent Pistes de l'onglet Findings pour ces cotes, non démontrées ici :
\begin{itemize}\small\setlength{\itemsep}{0pt}
\item Cote~23 : pour un anneau noethérien intègre $A$ quotient d'un anneau régulier, la clôture intégrale est finie sur $A$ si et seulement si les premiers de hauteur un où $A$ n'est pas régulier sont en nombre fini et que la normalisée de chaque $A_{\mathfrak p}$ correspondant est finie. Mathlib n'a pas les anneaux excellents ni les anneaux de Nagata.\piste{23-finite-normalisation-quotient-of-regular}
\item Cote~37 : pour une petite catégorie $C$ d'enveloppe de Karoubi $C'$, $\mathrm{Ind}(C)$ est un topos si et seulement si $C'$ a les limites inductives finies et que son foncteur vers les faisceaux pour la topologie finitaire les respecte. Mathlib n'a pas les conditions de SGA~4 pour que $\mathrm{Ind}(C)$ soit un topos.\piste{37-ind-topos-exercise}
\item Cote~37 : pour $B = k[X_0, X_1, \ldots]$ et $A = B[[T]]$, les modules $A_n/J_n$, $J_n = (X_0, X_1T, \ldots, X_nT^n)$, forment un système compatible de modules de présentation finie qui ne vient d'aucun $A$-module de présentation finie : la condition b) d'EGA~I 10.10.5 n'entraîne pas a) et c) sans hypothèse noethérienne. Accessible avec les polynômes et les séries formelles de mathlib ; l'anneau $B$ à une infinité de variables demande du soin.\piste{37-ega-i-10-10-5-counterexample}
\item Cote~90 : pour $G$ plat affine sur un schéma $S$ quelconque, d'algèbre réunion filtrante de sous-cogèbres localement libres, opérant sur $X$ avec un fibré ample $G$-linéarisé, tout sous-module quasi-cohérent d'un $G$-module quasi-cohérent est contenu dans un plus petit sous-module $G$-stable. Mathlib n'a pas les modules quasi-cohérents équivariants.\piste{90-equivariant-resolution-general-base}
\item Cote~112 : question : jusqu'où la correspondance de Dold--Kan subsiste hors de $\Delta$, restreinte aux préfaisceaux abéliens \q{lisses}. Mathlib n'a pas les catégories dérivées de complexes localement constants.\piste{112-linearization-beyond-the-simplex}
\item Cote~116 : énoncé sans démonstration : un topos recollé d'un sous-topos ouvert et de son fermé complémentaire est un topos de préfaisceaux exactement quand les deux le sont et que le foncteur de recollement commute aux limites projectives. Mathlib n'a pas le recollement des topos.\piste{116-presheaf-gluing-continuous-gluing-functor}
\item Cote~161-3 : $\mathrm{Top}(X \times Y) \to \mathrm{Top}(X) \times \mathrm{Top}(Y)$ est une équivalence sous une hypothèse plus faible que la compacité locale : toute partie localement fermée non vide de $X$ a un point qui y possède un voisinage relativement quasi-compact. Mathlib n'a pas les topos de faisceaux sur un produit.\piste{161-3-product-converse}
\item Cote~161-3 : pour une famille quelconque, $\mathrm{Top}(\prod X_i) \to \prod \mathrm{Top}(X_i)$ est une équivalence quand tous les facteurs sauf un au plus ont des systèmes fondamentaux de voisinages quasi-compacts, et tous sauf un au plus sont quasi-compacts. Mathlib n'a pas les produits infinis de topos de faisceaux.\piste{161-3-infinite-product-criterion}
\item Cote~161-3 : la spatialité du locale produit, démontrée par trois passages à la borne supérieure dans le treillis des sous-objets. Mathlib n'a pas les produits de locales.\piste{161-3-locale-spatiality-suprema-proof}
\item Cote~161-3 : le 2-produit de deux topos est l'intersection, dans le topos des préfaisceaux sur $C \times D$, des images réciproques des sous-topos $\tilde C \subset \hat C$ et $\tilde D \subset \hat D$. Mathlib n'a pas les sous-topos ni les 2-produits de topos.\piste{161-3-two-product-by-intersection}
\end{itemize}
